\documentclass[11pt]{article}
\usepackage[T1]{fontenc}
\usepackage{amsmath,amssymb,geometry}
\geometry{margin=1in}
\begin{document}
\begin{center}
{\Large Upper Laplacian kernels are not cohomology}\\[4pt]
Disproof of Conjecture 00000003713\\[4pt]
ziangni-sys --- October 8, 2026
\end{center}
\paragraph{The actual simplicial complex.}
Let $K$ be the filled two-simplex on vertices $0,1,2$: its faces
are every subset of these three vertices. Orient its three edges
as $01,02,12$ and its unique face as $012$. Work with rational
cochains and the standard coordinate dot products. Thus
$C^0(K;\mathbb Q)=\mathbb Q^3$, $C^1(K;\mathbb Q)=\mathbb Q^3$
and $C^2(K;\mathbb Q)=\mathbb Q$. The oriented incidence rule gives
\[
 d_0(a,b,c)=(b-a,c-a,c-b),\qquad d_1(x,y,z)=x-y+z.
\]
The second formula uses the actual oriented boundary
$\partial[012]=[12]-[02]+[01]$. Direct substitution shows
$d_1d_0=0$.

\paragraph{Genuine first cohomology.}
By definition,
\[
 H^1(K;\mathbb Q)=\ker d_1/\operatorname{im}d_0.
\]
If $(x,y,z)\in\ker d_1$, then $x-y+z=0$ and $z=y-x$.
The vertex potential $(0,x,y)$ consequently satisfies
$d_0(0,x,y)=(x,y,y-x)=(x,y,z)$. Hence
$\operatorname{im}d_0=\ker d_1$ and the actual quotient
$H^1(K;\mathbb Q)$ is the zero vector space.

\paragraph{Actual upper Laplacian.}
The degree-one upper Laplacian is $L_1^{\rm up}=d_1^*d_1$.
For the specified dot products, the adjoint is the transpose of
the face-edge incidence matrix $B=(1,-1,1)$, so
\[
 L_1^{\rm up}=B^{\mathsf T}B=
 \begin{pmatrix}1&-1&1\\-1&1&-1\\1&-1&1\end{pmatrix}.
\]
The nonzero cochain $v=(1,1,0)$ satisfies $Bv=0$, and therefore
$L_1^{\rm up}v=0$. Its kernel is consequently nonzero and cannot
be isomorphic to the zero first cohomology space. This contradicts
the claimed identification of the upper-Laplacian nullspace with
cohomology, disproving the conjecture.

\paragraph{Why the distinction matters.}
The full Hodge Laplacian includes the lower term
$d_0d_0^*$ as well. The standard harmonic-cohomology statement
concerns that full operator. The upper term alone also kills
nonzero coboundaries: our witness is explicitly
$v=d_0(0,1,1)$. Those coboundaries represent zero in cohomology,
but remain nonzero vectors in the upper kernel. The extra clause
about flow weights need not be resolved to disprove the conjunction.

\paragraph{Formal certificate.}
Lean constructs the full finite face family, certifies downward
closure and all face counts, and checks the actual three edges.
It defines the oriented coboundaries, proves the cochain identity,
forms the genuine quotient of cycles by boundaries and proves it
subsingleton. It certifies transpose adjointness for the dot
products, constructs the upper operator as the actual incidence
transpose product and verifies a nonzero kernel witness. The
final theorem rules out any linear equivalence between that kernel
and the actual cohomology quotient. Final audits use only standard
Lean and Mathlib axioms.
\end{document}
